\section{Relações Bem-fundadas e Set-like}
Adotemos, durante este capítulo, o termo "relação definida própria" como um símbolo de predicado binário da linguagem que está em uso, seja ele primitivo ou introduzido por definição e "relação definida" como uma relação definida própria ou uma relação conjunto.

Vamos definir a notação de relacão set-like

\begin{definition}
    Uma relação definida $ R $ é dita set-like se, e somente se,
    $$
        \pred_R(x) = \set{z : zRx}
    $$
    é um conjunto para todo $ x $.
\end{definition}

Para a parte de ser bem-fundada:

\begin{definition}
    Dado um conjunto $ X $, um elemento $x\in X$ é dito $R$-minimal se, e somente se, para todo $y\in X(y\centernot R x)$. Assim, $R$ é dita bem-fundada se
    $$
    \forall X(X\neq \emptyset \to \exists x\in A (y\in A \to y\centernot R x))
    $$
    Isto é, $R$ é bem-fundada se todo conjunto não-vazio possuir elemento $R$-minimal.
\end{definition}

\begin{remark}
    Notemos que quando $R$ acima é definida própria, as definições acima viram esquemas de definições, pois seria definido para cada símbolo de relação introduzido por definição. Omitiremos observações como esta ao decorrer da seção, mas o leitor será sábio em perceber que também se aplicará aos resultados e definições seguintes.
    
    Obviamente, toda relação conjunto $R$ é set-like, uma vez que $\pred_R(x) \subseteq \dom R$. 
\end{remark}

\begin{example}
    $\in$ é set-like, pois $pred_\in(x) = \{y: y\in x\}$ é um conjunto.
\end{example}
\begin{example}
    $pred_\preceq(1)$ seria o conjunto de todos os singletons, que não existe.   
\end{example}

\begin{example}
    $\subseteq$ é set-like se, e somente se, vale o Axioma da Potência. Nesse caso, temos que $pred_\subseteq(x) = \mathcal P(x)$.
\end{example}

\begin{example}
    = é set-like, pois $pred_=(x) = {x}$. Não, é, porém, bem-fundada, pois não possui um elemento =-minimal, uma vez que $x = x$. Analogamente, nenhuma relação reflexiva é bem-fundada.
\end{example}
\begin{example}
    $<$ ordenação para ordinais é set-like, pois $<$ é $\in|_{ON}$. Se $x$ é não vazio, todo elemento de $X$ que não é ordinal é minimal, senão, existe $\alpha = \min X$ e este será o $<$-minimal, pois $\forall \beta \in X, \alpha \leq \beta$, assim, $\neg(\beta < \alpha)$.
\end{example}


\begin{proposition}
    Sejam $R$ e $S$ relações definidas tais que para todo $x$ e para todo $y$, se $xRy$, então $xSy$, então:
    \begin{enumerate}[label=\roman*)]
        \item Se $S$ é set-like, então $R$ é set-like
        \item Se $S$ é bem-fundada, então $R$ é bem-fundada.
    \end{enumerate}
    Ademais, abrevia-se a hipótese dessa proposição como $R \subseteq S$ e dizemos que $S$ estende $R$.
\end{proposition}
\begin{proof}
     Para i), basta notar que, como $S$ é set-like, $\pred_S(x)$ é conjunto para todo $x$, assim, $\pred_R(x) = \set{y \in \pred_S(x): yRx}$ é conjunto pelo axioma da compreenção.

     Para ii), Consideremos $A$ um conjunto não vazio. Sejam $m\in A$ um elemento $S$-minimal de $A$, que sabemos que existe por hipótese. Então, $m$ será $R$-minimal, de fato, suponha que exista $y$ tal que $yRm$, teríamos, então que $ySm$ o que implica que $y\notin A$.
\end{proof}

Vamos ver, agora, como estender relações set-likes de forma que se transformem em relações transitivas.

\begin{definition}
    Seja $R$ uma relação definida. Para cada natural $n \geq 1$ e cada par $a,b$, um $R$-caminho de $n$ passos de $a$ até $b$ é uma sequência de elementos $(x_0,...,x_n)$ tal que $x_i R x_{i+1}$ para todo $i < n$ e que $x_0 = a$ e $x_n = b$.

    Assim, definamos $x R^*y$ se, e somente se, existe um caminho de $n \geq 1$, com $n$ natural, passos de $x$ a $y$. Ademais, dizemos que $R^*$ é o fecho transitivo, ou a transitivização, de $R$
    
\end{definition}

Provemos que $R^*$ é, de fato, transitivo e que extende $R$, justificando o nome dado na definição:

\begin{proposition}
    Se $R$ é uma relação definida, então $R^*$ é uma relação transitiva.
\end{proposition}
\begin{proof}
    É claro que $R^*$ estende $R$, uma vez que se $xRy$, então $xR^* y$, pelo caminho $(x,y)$. 
    
    Para transitividade, lembremos que queremos provar que, se $xR^*y$ e $yR^*z$, então $xR^*z$. Ora, suponha que a premissa é verdadeira, então temos dois $R$-caminhos, $(x, x_1,...,x_{n-1},y)$ e $(y,y_1,...,y_{m-1}, z)$ de $n$ e $m$ passos, respectivamente. Notemos que $(x,x_1,...,y,y_1,...z)$ é um $R$-caminho que liga $x$ a $z$, assim, concluímos que $xR^*z$, como queríamos.
\end{proof}

\begin{remark}
    Na verdade, $R^*$ é a menor relação transitiva que contém $R$. Vejamos, (PENDING)
\end{remark}

Mostremos que a transitivização herda a propriedade de ser set-like:

\begin{proposition}
    Se $R$ é uma relação definida set-like, então $R^*$ é uma relação definida set-like
\end{proposition}
\begin{proof}
    Seja $X$ um conjunto qualquer, provemos por indução que para todo $n \geq 1$ natural existe um único $y$ tal que para todo $z \in y$, $z\in y \iff zR^*x$ com um caminho de $n$ passos. Chamemos $y$ de $\pred_R^n(x)$.

    Para $n=1$, segue de que $R$ é set-like. Suponha, agora, que vale para $n \geq 1$, pelo Esquema da Substituição, o conjunto $A = \set{\pred_R(y) : y \in \pred_R^n(x)}$ é bem definido, pois $\pred_R(y)$ é único para todo $y\in \pred_R^n(x)$. Sabendo disso, definamos $\pred_R^{n+1}(x) = \bigcup A$. Mostremos que vale o que queremos, isto é, que é o conjunto de todos os $y$ tais que $yR^*x$ por um $R$-caminho de $n + 1$ passos.
    
    Claramente todos elementos $z \in \pred_R^{n+1}(x)$ tem $R$-caminho de $n + 1$ passos para $x$, pois basta pegar $y$ tal que $z\in pred_R^n(y) \in A$, tomar o $R$-caminho de $y$ a $x$ e concatenar $z$ no início. 
    
    Reciprocamente, vamos mostrar que todos $z$ tais que $zR^*x$ por um $R$-caminho de $n + 1$ passos estão em $\pred_R^{n+1}(x)$. Seja $(x_0, x_1,...,x_n,x)$ um caminho de $n + 1$ passos, sabemos que $x_1 \in \pred_R^n(x)$ e que $x_0R x_1$, logo, $x_0 \in \pred_R(x_1)$ o que implica que $x\in pred_R^{n+1}(x)$, concluíndo a indução. 

    Com isso e finalmente, consideremos $Z = \bigcup_{0 < n \in \omega}\pred_R^n(x)$, por Substuição e Axioma da União, é um conjunto e temos que $zR^*x \iff z \in Z$, assim, provamos que $R^*$ é set-like, como queríamos.
\end{proof}

\begin{definition}
    Seja $R$ uma relação definida e $A$ uma classe. A restrição de $R$ a $A$, denotada por $R_A$, é definida por:

    $$
    \forall x\forall y(xR_A y \longlif (x\in A \land y\in A \land xRy))
    $$

    Quando $A$ é um conjunto, $R_A = \set((x,y)\in A : xRy)$. Dizemos que $R$ é bem-fundada em $A$ se $R_A$ é bem-fundada. 
\end{definition}

Convencionemos $R^*_A = (R_A)^*$ e não $(R^*)_A$. Como um check de sanidade, provemos:

\begin{lemma}
    Seja $R$ uma relação definida e $A$ uma classe. Então, $R$ é bem-fundada em $A$ se, e somente se, para todo $X \subseteq A$ existe um $R$-minimal de $X$.
\end{lemma}
\begin{proof}
    
\end{proof}